
\phantomsection
\section*{Глава 3. Программная реализация веб-приложения}
\addcontentsline{toc}{section}{Глава 3. Программная реализация веб-приложения}
\stepcounter{section}

\subsection{Структура программного решения}

Программная реализация веб-приложения для частной медицинской клиники выполнена на объектно-ориентированном языке программирования Python с использованием микрофреймворка Flask. Выбор Flask обусловлен его легковесностью, модульностью и простотой интеграции со сторонними библиотеками для работы с базами данных (Flask-SQLAlchemy) и организации сессий пользователей (Flask-Login).

Проект имеет следующую структуру директорий и файлов:
\begin{enumerate}
    \item \texttt{app.py} - основной файл приложения, содержащий инициализацию Flask-сервера, настройку обработчиков маршрутов (endpoints), логику обработки запросов и авторизации.
    \item \texttt{models.py} - модуль описания структуры базы данных на уровне объектно-реляционного отображения (ORM SQLAlchemy).
    \item \texttt{config.py} - конфигурационный класс приложения, определяющий настройки безопасности (секретный ключ сессии), путь к базе данных SQLite/PostgreSQL, параметры загрузки документов и лимиты на размер файлов.
    \item \texttt{seed.py} - скрипт для первоначального создания таблиц базы данных и наполнения их демонстрационными данными (учетными записями тестовых пользователей различных ролей, списком услуг и расписанием врачей).
    \item \texttt{requirements.txt} - перечень внешних зависимостей проекта и версий библиотек для сборки окружения.
    \item \texttt{templates/} - каталог HTML-шаблонов страниц с разметкой на базе шаблонизатора Jinja2.
    \item \texttt{static/} - папка для статических ресурсов приложения, включающая стили CSS, скрипты JavaScript, иконки Bootstrap Icons и подкаталог \texttt{uploads/} для загрузки медицинских заключений и файлов пациентов.
\end{enumerate}

Конфигурационный файл \texttt{config.py} обеспечивает адаптивность приложения к среде развертывания: при наличии в переменных окружения строки подключения \texttt{DATABASE\_URL} приложение подключается к внешней СУБД PostgreSQL (промышленный режим), иначе используется локальный файл базы данных \texttt{clinic.db} через встроенную поддержку SQLite (режим разработки).

\subsection{Модели данных и ORM-маппинг}

Для управления персистентным состоянием и выполнения запросов к СУБД в приложении применена технология ORM (Object-Relational Mapping) на базе Flask-SQLAlchemy. Описание логических сущностей и связей выполнено в виде Python-классов, наследующих базовый класс декларативного описания \texttt{db.Model}.

Основные классы моделей, реализованные в модуле \texttt{models.py}:
\begin{enumerate}
    \item \texttt{User} - учётная запись пользователя. Содержит поля \texttt{id}, \texttt{email}, \texttt{password\_hash} (хэш пароля на базе алгоритма scrypt) и \texttt{role} (роль: \texttt{'patient'}, \texttt{'doctor'} или \texttt{'admin'}). Реализует методы \texttt{set\_password} для хэширования и \texttt{check\_password} для проверки введённого пароля.
    \item \texttt{Patient} - профиль пациента. Связан отношением 1:1 с учётной записью \texttt{User} с использованием каскадного удаления \texttt{cascade="all, delete-orphan"}. Содержит персональные данные (ФИО, дата рождения, телефон).
    \item \texttt{Doctor} - профиль врача. Связан отношением 1:1 с учётной записью \texttt{User}. Дополнительно содержит поля специализации (ссылка на \texttt{Specialization}) и номера кабинета приёма \texttt{cabinet}.
    \item \texttt{Specialization} - справочник врачебных специальностей (например, терапевт, кардиолог, невролог).
    \item \texttt{Service} - перечень предоставляемых услуг с указанием названия, описания, цены и длительности в минутах.
    \item \texttt{ScheduleSlot} - интервалы расписания приёма. Содержит ссылку на \texttt{Doctor}, дату, время начала и окончания интервала, а также статус занятости слота (\texttt{'free'}, \texttt{'reserved'}).
    \item \texttt{Appointment} - запись на приём. Связывает пациента (\texttt{Patient}), врача (\texttt{Doctor}), услугу (\texttt{Service}) и временной интервал расписания (\texttt{ScheduleSlot}). Поле \texttt{slot\_id} снабжено ограничением уникальности \texttt{unique=True}, предотвращающим повторное бронирование одного слота.
    \item \texttt{MedicalRecord} - электронная медицинская карта. Хранит группу крови, резус-фактор, перечень аллергенов и хронических болезней пациента. Связана 1:1 с пациентом.
    \item \texttt{Visit} - протокол осмотра (посещения). Содержит жалобы, диагноз и рекомендации врача. Создаётся для конкретной записи \texttt{Appointment} после выполнения приёма.
    \item \texttt{Document} - прикрепленные файлы исследований (картинки, документы PDF). Связан с пациентом и (опционально) с конкретным посещением.
\end{enumerate}

Связи между таблицами настроены с использованием внешних ключей (\texttt{db.ForeignKey}) с поддержкой ограничений целостности БД (\texttt{ondelete='CASCADE'} и \texttt{ondelete='SET NULL'}), что обеспечивает корректную работу транзакций СУБД.

Пример реализации структуры моделей данных на языке Python с использованием библиотеки Flask-SQLAlchemy представлен в листинге \ref{lst:models_code}.

\begin{lstlisting}[language=Python, caption={Реализация моделей данных в clinic\_app/models.py}, label={lst:models_code}]
from flask_sqlalchemy import SQLAlchemy
from flask_login import UserMixin
from werkzeug.security import generate_password_hash, check_password_hash
from datetime import datetime

db = SQLAlchemy()

class User(db.Model, UserMixin):
    __tablename__ = 'users'
    id = db.Column(db.Integer, primary_key=True)
    email = db.Column(db.String(120), unique=True, nullable=False)
    password_hash = db.Column(db.String(256), nullable=False)
    role = db.Column(db.String(20), nullable=False) # 'patient', 'doctor', 'admin'
    created_at = db.Column(db.DateTime, default=datetime.utcnow)
    
    patient = db.relationship('Patient', backref='user', uselist=False, cascade="all, delete-orphan")
    doctor = db.relationship('Doctor', backref='user', uselist=False, cascade="all, delete-orphan")
    
    def set_password(self, password):
        self.password_hash = generate_password_hash(password)
        
    def check_password(self, password):
        return check_password_hash(self.password_hash, password)

class Patient(db.Model):
    __tablename__ = 'patients'
    id = db.Column(db.Integer, primary_key=True)
    user_id = db.Column(db.Integer, db.ForeignKey('users.id', ondelete='CASCADE'), unique=True, nullable=False)
    full_name = db.Column(db.String(150), nullable=False)
    birth_date = db.Column(db.Date, nullable=False)
    phone = db.Column(db.String(20), nullable=False)
    
    appointments = db.relationship('Appointment', backref='patient', lazy=True, cascade="all, delete-orphan")
    medical_record = db.relationship('MedicalRecord', backref='patient', uselist=False, cascade="all, delete-orphan")

class Appointment(db.Model):
    __tablename__ = 'appointments'
    id = db.Column(db.Integer, primary_key=True)
    patient_id = db.Column(db.Integer, db.ForeignKey('patients.id', ondelete='CASCADE'), nullable=False)
    doctor_id = db.Column(db.Integer, db.ForeignKey('doctors.id'), nullable=False)
    service_id = db.Column(db.Integer, db.ForeignKey('services.id'), nullable=False)
    slot_id = db.Column(db.Integer, db.ForeignKey('schedule_slots.id'), unique=True, nullable=False)
    status = db.Column(db.String(20), nullable=False, default='created')
    created_at = db.Column(db.DateTime, default=datetime.utcnow)
    visit = db.relationship('Visit', backref='appointment', uselist=False, cascade="all, delete-orphan")
\end{lstlisting}

\subsection{Логика маршрутизации и бизнес-логика}

Маршрутизация запросов и бизнес-логика реализованы в виде контроллеров Flask в файле \texttt{app.py}. Доступ к защищённым страницам регулируется расширением \texttt{Flask-Login}. Ограничение доступа к страницам личных кабинетов реализовано с помощью декоратора \texttt{@login\_required} и проверок роли текущего авторизованного пользователя \texttt{current\_user.role}.

Основные функциональные маршруты приложения:
\begin{enumerate}
    \item \textbf{Публичные маршруты}:
    \begin{itemize}
        \item \texttt{/} (функция \texttt{index}) - отображение главной страницы с направлениями работы и списком услуг клиники;
        \item \texttt{/auth} (функция \texttt{auth\_page}) - страница аутентификации и регистрации;
        \item \texttt{/login} и \texttt{/register} (POST-обработчики) - выполняют валидацию входных данных, проверку уникальности email, авторизацию сессии с помощью \texttt{login\_user} и перенаправление пользователя в личный кабинет в зависимости от его роли.
    \end{itemize}

    \item \textbf{Личный кабинет пациента}:
    \begin{itemize}
        \item \texttt{/cabinet} (функция \texttt{patient\_cabinet}) - собирает информацию о профиле пациента, его электронной медицинской карте, а также списки будущих записей на приём и архив прошедших визитов с рекомендациями врачей;
        \item \texttt{/booking} (функция \texttt{booking\_page}) - страница онлайн-записи на приём. Передаёт в шаблон списки направлений, врачей и услуг;
        \item \texttt{/booking/get\_slots} - API-маршрут (возвращает JSON), используемый динамическим JavaScript на клиенте для получения свободных слотов расписания конкретного врача на выбранную дату;
        \item \texttt{/booking/create} (POST-запрос) - создание записи на приём. Бронирует выбранный слот (переводит его статус в \texttt{'reserved'}), создаёт объект \texttt{Appointment} и перенаправляет в кабинет.
    \end{itemize}

    \item \textbf{Личный кабинет врача}:
    \begin{itemize}
        \item \texttt{/doctor} (функция \texttt{doctor\_cabinet}) - отображает расписание приёмов врача на выбранную дату (по умолчанию на текущий день), позволяя фильтровать записи;
        \item \texttt{/doctor/visit/<int:appt\_id>} - страница проведения приёма пациента. Врач вносит текстовые жалобы, диагноз и рекомендации. По завершении создаётся запись \texttt{Visit}, статус приёма обновляется на \texttt{'completed'}, а слот расписания закрывается;
        \item \texttt{/doctor/upload\_document} - обработчик загрузки медицинских файлов. Выполняет проверку расширения файла, сохраняет его в папку \texttt{static/uploads} с безопасным именем и связывает запись \texttt{Document} с пациентом.
    \end{itemize}

    \item \textbf{Панель администратора}:
    \begin{itemize}
        \item \texttt{/admin} (функция \texttt{admin\_dashboard}) - сводная панель со статистикой работы клиники;
        \item \texttt{/admin/doctors} и \texttt{/admin/services} - управление справочниками медицинского персонала и услуг (добавление, редактирование, удаление записей);
        \item \texttt{/admin/appointments} - общий реестр записей на приём с возможностью подтверждения (\texttt{confirm}) или отмены (\texttt{cancel}).
    \end{itemize}
\end{enumerate}

При обработке форм реализован механизм всплывающих уведомлений \texttt{flash()}, информирующий пользователя об успешном завершении операции или возникших ошибках валидации данных.

\newpage

Пример реализации обработчика маршрута для подтверждения записи пациента на приём с контролем занятости выбранного слота расписания и транзакционной логикой СУБД представлен в листинге \ref{lst:routing_code}.

\begin{lstlisting}[language=Python, caption={Реализация маршрута подтверждения записи в clinic\_app/app.py}, label={lst:routing_code}]
@app.route('/confirm_booking', methods=['POST'])
@login_required
def confirm_booking():
    slot_id = request.form.get('slot_id')
    doctor_id = request.form.get('doctor_id')
    service_id = request.form.get('service_id')
    
    slot = ScheduleSlot.query.get(slot_id)
    if not slot or slot.status != 'free':
        flash('Выбранный временной слот уже занят.', 'error')
        return redirect(url_for('booking_page'))
        
    slot_datetime = datetime.combine(slot.date, slot.start_time)
    if slot_datetime < datetime.now():
        flash('Выбранное время приема уже прошло.', 'error')
        return redirect(url_for('booking_page'))
        
    try:
        appt = Appointment(
            patient_id=current_user.patient.id,
            doctor_id=int(doctor_id),
            service_id=int(service_id),
            slot_id=int(slot_id),
            status='confirmed'
        )
        db.session.add(appt)
        slot.status = 'booked'
        db.session.commit()
        
        flash('Вы успешно записались на приём!', 'success')
        return redirect(url_for('patient_cabinet'))
        
    except Exception as e:
        db.session.rollback()
        flash(f'Ошибка при подтверждении записи: {str(e)}', 'error')
        return redirect(url_for('booking_page'))
\end{lstlisting}

\subsection{Интерфейс пользователя и экраны приложения}

Пользовательский интерфейс спроектирован в соответствии с современными требованиями к веб-ресурсам и построен на адаптивной CSS-сетке фреймворка Bootstrap. Дизайн выполнен в приятной светлой теме с использованием корпоративных синих, бирюзовых и серых оттенков, качественной типографики и иконок Bootstrap Icons.

В ходе тестирования с помощью автоматического агента были зафиксированы скриншоты работающего веб-приложения для всех ключевых сценариев использования. Для отделения изображений от белого фона страниц и придания им завершенного вида все снимки экранов снабжены тонкой рамкой и представлены в виде отдельных рисунков в крупном формате.

\setlength{\fboxsep}{0pt}%
\setlength{\fboxrule}{0.5pt}%

Главная страница приложения представлена на рисунке \ref{fig:screenshot_main}. На ней размещены общая информация о клинике, направления работы со справочными иконками, карточки доступных услуг с указанием стоимости и кнопка быстрого перехода к онлайн-записи.

\begin{figure}[ht!]
\centering
\fbox{\includegraphics[width=0.90\textwidth,height=0.40\textheight,keepaspectratio]{screenshot_main}}
\caption{Скриншот главной страницы веб-приложения}
\label{fig:screenshot_main}
\end{figure}

Страница входа и регистрации представлена на рисунке \ref{fig:screenshot_auth}. Она содержит две формы: форму авторизации для зарегистрированных пациентов и сотрудников (слева) и форму регистрации новых пациентов (справа), снабженную обязательным предупреждением о согласии на обработку персональных данных в соответствии с ФЗ-152.

\begin{figure}[ht!]
\centering
\fbox{\includegraphics[width=0.90\textwidth,height=0.40\textheight,keepaspectratio]{screenshot_auth}}
\caption{Скриншот страницы авторизации и регистрации}
\label{fig:screenshot_auth}
\end{figure}

Личный кабинет пациента представлен на рисунке \ref{fig:screenshot_patient}. В верхней части отображаются данные электронной медицинской карты пациента (группа крови, резус-фактор, список аллергенов и противопоказаний). Ниже расположены таблицы с предстоящими приёмами (с возможностью отмены) и архивом прошедших посещений врача.

\begin{figure}[ht!]
\centering
\fbox{\includegraphics[width=0.90\textwidth,height=0.40\textheight,keepaspectratio]{screenshot_patient}}
\caption{Скриншот личного кабинета пациента}
\label{fig:screenshot_patient}
\end{figure}

Форма онлайн-записи на приём представлена на рисунке \ref{fig:screenshot_booking}. Она реализует пошаговый интерфейс выбора направления, врача и услуги. Сетка свободных слотов расписания подгружается динамически на основе запросов к СУБД с использованием асинхронного JavaScript.

\begin{figure}[ht!]
\centering
\fbox{\includegraphics[width=0.90\textwidth,height=0.40\textheight,keepaspectratio]{screenshot_booking}}
\caption{Скриншот процесса онлайн-записи на приём}
\label{fig:screenshot_booking}
\end{figure}

Личный кабинет врача (рабочее место доктора) представлен на рисунке \ref{fig:screenshot_doctor}. В левой части экрана отображается сетка записей на выбранный день. При выборе конкретной записи справа открывается интерактивная медицинская форма приёма для фиксации жалоб, диагноза, рекомендаций и прикрепления файлов результатов исследований.

\begin{figure}[ht!]
\centering
\fbox{\includegraphics[width=0.90\textwidth,height=0.40\textheight,keepaspectratio]{screenshot_doctor}}
\caption{Скриншот рабочего места врача при заполнении приёма}
\label{fig:screenshot_doctor}
\end{figure}

\clearpage

Панель управления администратора клиники представлена на рисунке \ref{fig:screenshot_admin}. Она предоставляет полноценный CRUD-интерфейс для добавления медицинских услуг, создания учетных записей врачей, формирования сетки расписания и контроля записей пациентов.

\begin{figure}[ht!]
\centering
\fbox{\includegraphics[width=0.90\textwidth,height=0.40\textheight,keepaspectratio]{screenshot_admin}}
\caption{Скриншот панели управления администратора клиники}
\label{fig:screenshot_admin}
\end{figure}

Таким образом, разработанное программное обеспечение полностью решает поставленные задачи по автоматизации частной клиники, предоставляя удобные интерфейсы для всех категорий пользователей и гарантируя надежную обработку информации на бэкенде.
